\documentclass[
sansserifheadings, % Für Nutzung von Serifen-Überschriften ein Prozentzeichen am Anfang setzen
acrosingle, % Nimmt nur Abkürzungen ins Inhaltsverzeichnis auf, die mehr als einmal verwendet wurden
%acroindex, % Nimmt die verwendeten Abkürzungen als Einträge ins Stichwortverzeichnis auf
%englisch, % Stellt das Dokument auf Englisch um
%abbreviationstoc, % Nimmt einen Einträg für das Abkürzungsverzeichnis ins Inhaltsverzeichnis mit auf
%twoside % Stellt das Dokument auf eine zweiseitige Ausgabe um -- gegebenenfalls müssen in diesem Fall die Seitenränder etwas angepasst werden
]{politikunibremen}

% Einbindung der Daten (nicht nur) für das Deckblatt %

\input{Deckblatt/DatenAutorIn.tex}      
\input{Deckblatt/DatenDeckblatt.tex}
\input{CD-Cover/CD-Etikett.tex}

% Daten für das Abkürzungsverzeichnis %

\input{Acronyms.tex} 

% Platz für weitere Pakete %

\usepackage{pgfplots} \pgfplotsset{compat=1.9} % Für das Zeichnen von Graphen auf der Grundlage von csv-Dateien

            %%% Hier beginnt das Dokument %%%
  
\begin{document}  

% Einfügung des Deckblatts %

\input{Deckblatt/Deckblatt.tex}

\Danksagung{
 
Ich möchte vor allem meinem guten Freund und Zweigutachter John H. Watson danken, der mir auf fachlicher wie auch auf menschlicher Ebene stets weiterhelfen konnte und wertvolle Impulse für die wissenschaftliche Arbeit gab. 
Mein Dank gilt weiterhin meiner Vermieterin Mrs. Hudson, die die Arbeit mehrfach mit einem kritischen Blick durchsah und mir bei unzähligen Tassen Tee Feedback aus der Sicht einer Leserin gab, die unvoreingenommen mit dem Thema der Arbeit konfrontiert wurde.

}

\Widmung{Für Mycroft.} % Seite mit Widmung

% Seite mit dem Inhaltsverzeichnis

\Inhaltsverzeichnis % Seite mit dem Inhaltsverzeichnis

\Abkuerzungsverzeichnis % Seite mit Abkürzungsverzeichnis

                          %%% Hier beginnt eure Arbeit! %%%      

\nocite{*} % Stellt auch nicht zitierte Literatur (zur Veranschaulichung des Zitierstils) im Literaturverzeichnis dar

\section{Einleitung}
        
Dies ist der erste Satz der Einleitung!\footnote{Dies ist eine supertolle Fußnote!} Mit \ac{LaTeX}\index{LaTeX} lassen sich sehr schöne Dokumente erstellen, auch wenn es teilweise etwas kompliziert sein kann, wenn eins z. B. Deckblätter bearbeiten möchte -- ich spreche aus Erfahrung (vgl. hierzu auch \citealt{Weller.1994, Moriarty.2011}). Prof. Dr. James \citet{Moriarty.2011b} -- seines Zeichens Leiter des kriminologischen Instituts\index{Universität Bremen!Kriminologisches Institut} der Universität Bremen\index{Universität Bremen} -- hat zum Thema der Arbeit einen ganzen Sammelband herausgebracht, wie beispielhaft erwähnt sei. In seiner bahnbrechenden Publikation \enquote{The Role of Supervillains in Modern Democracy} arbeitet \citet{Moriarty.2011} die Kernkontroversen des Gegenstandes der vorliegenden Arbeit wie folgt heraus:   

\begin{quote}
\blindtext Abschließend lässt sich festhalten: Das Schreiben von Blindtext liegt einfach in meiner \ac{DNA}\index{DNA}.
\end{quote} 
\blindtext

\subsection{Theoretischer Rahmen} 

Hier könnten beispielsweise eure fundierten theoretischen Überlegungen stehen.

\subsubsection{Hypothese 1}

Neben Subsections sind -- wie ihr seht -- auch Subsubsections möglich.

\begin{figure}[ht]
\centering
\includegraphics[width=0.5\textwidth]{Bilder/BeispielBild.jpg}
\caption[Ein Beispielbild]{Ein Beispielbild (CC BY-SA 2.0 Mike Quinn)}
\end{figure}  

\subsubsection{Hypothese 2}   

Aller guten Dinge sind ... zwei? Dies wurde bereits von \citet[121-232]{Holmes.2014} gezeigt.

\begin{table}[ht]
\centering
\begin{tabular}{|l|c|c|r|}  
\hline
\textbf{Ein} & \textbf{wirklich} & \textbf{grandioses} & \textbf{Beispiel} \\ \hline
für & eine & hervorragende & Tabelle! \\ \hline
\end{tabular}
\caption[Eine Beispieltabelle]{Eine Beispieltabelle.}
\end{table}

  
\subsection{Methodische Vorgehensweise} 
Hier könnte zum Beispiel stehen, wie ihr vorhabt zu versuchen, qualitativ hochwertige Forschung durchzuführen \citep{Watson.2013}!

\begin{itemize}
\item Ein Beispiel
\item für Aufzählungen
\item in \ac{LaTeX}\index{LaTeX!Aufzählungen}
\end{itemize}

\blindtext
  
\section{Hauptteil}

Hier beginnt der Hauptteil eurer Arbeit (siehe hierzu auch \citealt{Grimberg.2012}). Viel Spaß!

\begin{enumerate}
\item Ein Beispiel
\item für nummerierte Aufzählungen
\item in \ac{LaTeX}\index{LaTeX!Nummerierte Aufzählungen}
\end{enumerate}

\blindtext

\section{Noch mehr Hauptteil}

Wissenschaft kann eine langwierige Angelegenheit sein \citep[543ff]{Moriarty.2011}. Es folgen Beispielgleichungen: 

\begin{align} 
x_1+x_2+...+x_n &= -a_{n-1} \label{Summe} \\ 
x_1x_2+x_1x_3+...+x_{n-1}x_n &= a_{n-3}\\ 
x_1x_2 \cdot ... \cdot x_{n-1}x_n &= (-1)^na_0 
\end{align} 

\blindtext

\section{Schlussbetrachtung}

\Blindtext

Das war's!


% Verzeichnisse -- nicht Benötigtes entfernen
      
\Literaturverzeichnis % Erstellt das Literaturverzeichnis. Die Daten sind in der bibo.bib-Datei. Der Bibliografie-Stil kann bei Bedarf in der politikunibremen.cls-Datei geändert werden
  
\Tabellenverzeichnis % Erstellt ein Tabellenverzeichnis  
  
\Abbildungsverzeichnis % Erstellt ein Abbildungsverzeichnisses

\Stichwortverzeichnis % Erstellt einen Index mit den zuvor festgelegten Stichwörtern. Erscheint nur, wenn mittels \index{Indexeintrag!Indexuntereintrag} Indexeinträge definiert wurden

% Hier kann der Anhang eingefügt werden %

\sectionanhang{Tabellen}
          
\begin{table}[H]
\centering
    \begin{tabular}{ | l | l | l | p{5cm} |}
    \hline
    Day & Min Temp & Max Temp & Summary \\ \hline
    Monday & 11C & 22C & A clear day with lots of sunshine.  
    However, the strong breeze will bring down the temperatures. \\ \hline
    Tuesday & 9C & 19C & Cloudy with rain, across many northern regions. Clear spells 
    across most of Scotland and Northern Ireland, 
    but rain reaching the far northwest. \\ \hline
    Wednesday & 10C & 21C & Rain will still linger for the morning. 
    Conditions will improve by early afternoon and continue 
    throughout the evening. \\
    \hline
    \end{tabular}
\caption{Noch eine Beispieltabelle.}
\label{Beispiel1}
\end{table}

\begin{table}[H]
\centering
    \begin{tabular}{ | l | l | l | p{5cm} |}
    \hline
    Day & Min Temp & Max Temp & Summary \\ \hline
    Monday & 11C & 22C & A clear day with lots of sunshine.  
    However, the strong breeze will bring down the temperatures. \\ \hline
    Tuesday & 9C & 19C & Cloudy with rain, across many northern regions. Clear spells 
    across most of Scotland and Northern Ireland, 
    but rain reaching the far northwest. \\ \hline
    Wednesday & 10C & 21C & Rain will still linger for the morning. 
    Conditions will improve by early afternoon and continue 
    throughout the evening. \\
    \hline
    \end{tabular}
\caption{Und noch eine Beispieltabelle.}
\label{Beispiel2}
\end{table}

\sectionanhang{Abbildungen}

\begin{figure}[H]
\centering
\includegraphics[width=0.5\textwidth]{Bilder/BeispielBild.jpg}
\caption[Noch ein Beispielbild]{Noch ein Beispielbild (CC BY-SA 2.0 Mike Quinn)}
\end{figure} 

\begin{figure}[H]
\centering
\includegraphics[width=0.5\textwidth]{Bilder/BeispielBild.jpg}
\caption[Und noch ein Beispielbild]{Und noch ein Beispielbild (CC BY-SA 2.0 Mike Quinn)}
\end{figure} 

% Schriftliche Versicherung als PDF - hier entweder \OptionEins, \OptionZwei oder \OptionDrei einfügen

\SchriftlicheVersicherungAbschlussarbeit{\OptionEins}

% Erstellung von CD-Etiketten, die auf die abgegebene CD geklebt werden können %

\CDEtikett{\TitelDerArbeit}{black} % Bei Formatierungsproblemen statt \TitelDerArbeit den Titel manuell eingeben. Die Umrisse verschwinden, wenn man statt "black" "white" (ohne Anführungszeichen) in die zweiten geschwungenen Klammern einträgt
  
\end{document} 
